%# -*- coding:utf-8 -*-
%% start of file `template_en.tex'.
%% Copyright 2006-1008 Xavier Danaux (xdanaux@gmail.com).
%
% This work may be distributed and/or modified under the
% conditions of the LaTeX Project Public License version 1.3c,
% available at http://www.latex-project.org/lppl/.


\documentclass[11pt,a4paper]{moderncv}

\usepackage{fontspec,xunicode}
\usepackage[slantfont,boldfont]{xeCJK}
\usepackage{xcolor}  % replace by the encoding you are using

\usepackage{ifplatform} % 判断平台


% \setmainfont{Tahoma}
\setmainfont{Times New Roman}  % 缺省英文字体.serif是有衬线字体sans serif无衬线字体

\ifwindows
  % Windows 下使用的代码
  \setCJKmainfont[ItalicFont=KaiTi, BoldFont=SimHei]{SimSun}  % 衬线字体 缺省中文字体为
  \setCJKsansfont{SimSun}
  \setmonofont{FangSong}                 % 等宽字体: FangSong, Consolas
\fi

\ifmacosx
  % macOS 下使用的代码
  \setCJKmainfont[ItalicFont={Kai}, BoldFont={Hei}]{STSong}  % 衬线字体 缺省中文字体为
  \setCJKsansfont{STSong}
  \setCJKmonofont{STFangsong}  % 中文等宽字体
\fi


%-----------------------xeCJK下设置中文字体------------------------------%
\setCJKfamilyfont{song}{SimSun}  % 宋体 song
\newcommand{\song}{\CJKfamily{song}}
\setCJKfamilyfont{fs}{FangSong_GB2312}  % 仿宋2312 fs
\newcommand{\fs}{\CJKfamily{fs}}
\setCJKfamilyfont{yh}{Microsoft YaHei}  % 微软雅黑 yh
\newcommand{\yh}{\CJKfamily{yh}}
\setCJKfamilyfont{hei}{SimHei}  % 黑体  hei
\newcommand{\hei}{\CJKfamily{hei}}
\setCJKfamilyfont{hwxh}{STXihei}  % 华文细黑  hwxh
\newcommand{\hwxh}{\CJKfamily{hwxh}}
\setCJKfamilyfont{asong}{Adobe Song Std}  % Adobe 宋体  asong
\newcommand{\asong}{\CJKfamily{asong}}
\setCJKfamilyfont{ahei}{Adobe Heiti Std}  % Adobe 黑体  ahei
\newcommand{\ahei}{\CJKfamily{ahei}}
\setCJKfamilyfont{akai}{Adobe Kaiti Std}  % Adobe 楷体  akai
\newcommand{\akai}{\CJKfamily{akai}}


%------------------------------设置字体大小------------------------%
\newcommand{\chuhao}{\fontsize{42pt}{\baselineskip}\selectfont}  % 初号
\newcommand{\xiaochuhao}{\fontsize{36pt}{\baselineskip}\selectfont}  % 小初号
\newcommand{\yihao}{\fontsize{28pt}{\baselineskip}\selectfont}  % 一号
\newcommand{\erhao}{\fontsize{21pt}{\baselineskip}\selectfont}  % 二号
\newcommand{\xiaoerhao}{\fontsize{18pt}{\baselineskip}\selectfont}  % 小二号
\newcommand{\sanhao}{\fontsize{15.75pt}{\baselineskip}\selectfont}  % 三号
\newcommand{\sihao}{\fontsize{14pt}{\baselineskip}\selectfont}  % 四号
\newcommand{\xiaosihao}{\fontsize{12pt}{\baselineskip}\selectfont}  % 小四号
\newcommand{\wuhao}{\fontsize{10.5pt}{\baselineskip}\selectfont}  % 五号
\newcommand{\subwuhao}{\fontsize{10pt}{\baselineskip}\selectfont}  % 次五号
\newcommand{\xiaowuhao}{\fontsize{9pt}{\baselineskip}\selectfont}  % 小五号
\newcommand{\liuhao}{\fontsize{7.875pt}{\baselineskip}\selectfont}  % 六号
\newcommand{\qihao}{\fontsize{5.25pt}{\baselineskip}\selectfont}  % 七号


% \usepackage{fontawesome}
% \setCJKmainfont[BoldFont={WenQuanYi Micro Hei/Bold}]{WenQuanYi Micro Hei}
% \defaultfontfeatures{Mapping=tex-text}
% \XeTeXlinebreaklocale "zh"
% \XeTeXlinebreakskip = 0pt plus 1pt minus 0.1pt
% moderncv themes
\moderncvtheme[blue]{classic}  % optional argument are 'blue' (default), 'orange', 'red', 'green', 'grey' and 'roman' (for roman fonts, instead of sans serif fonts)
% \moderncvtheme[green]{classic}  % idem
% \moderncvtheme[blue,roman]{hht}
% character encoding



% adjust the page margins
\usepackage[scale=0.88, top=6mm, bottom=6mm]{geometry}
\setlength{\hintscolumnwidth}{2.5cm}
\AtBeginDocument{\setlength{\maketitlenamewidth}{5cm}}
\AtBeginDocument{\recomputelengths}

% compact itemize spacing
\usepackage{enumitem}
\setlist[itemize]{nosep, leftmargin=1.5em, topsep=1pt, partopsep=0pt, parsep=0pt, itemsep=1pt}

% personal data
\firstname{詹}
\familyname{勇}
\title{个人简历}
\address{浙江理工大学}{}
\mobile{18856603434}
\email{3446914219@qq.com}
\social[github]{daetz-coder}
\extrainfo{WeChat: daetz\textunderscore{}zy}

\photo[48pt]{avatar.png}  % '64pt' is the height the picture must be resized to and 'picture' is the name of the picture file; optional, remove the line if not wanted
% \quote{China\TeX 您的LaTeX乐园，TeX\&\LaTeX 王国}  % optional, remove the line if not wante

% \nopagenumbers{}  % uncomment to suppress automatic page numbering for CVs longer than one page


%----------------------------------------------------------------------------------
%            content
%----------------------------------------------------------------------------------
\begin{document}
\small
\maketitle
\vspace*{-10mm}

\section{求职意向}
\cvline{目标岗位}{Agent 开发工程师}
\cvline{技术方向}{LLM Agent 系统架构、RAG 检索增强生成、AI 应用全栈开发}
\cvline{到岗时间}{随时到岗，可全职实习}

\section{教育经历}
\cventry{2024.09--至今}{硕士}{浙江理工大学}{计算机技术}{研究方向：AI Agent 与智能系统}{}
\cvlistitem{一等优秀奖学金（1次）、二等优秀奖学金（1次）}
\cventry{2020.09--2024.06}{本科}{嘉兴大学}{软件工程}{校级优秀毕业生}{}
\cvlistitem{浙江省政府奖学金（1次）、二等学习优秀奖学金（1次）、三等学习优秀奖学金（5次）}

\section{项目经历}
\cventry{2025.05--至今}{Wiki-Agent：AI 驱动的智能知识库系统}{独立开发者}{}{}{%
  基于 LangGraph 的 AI Agent 系统，以对话形式管理知识库，支持 Human-in-the-Loop 机制。\\
  \textbf{技术栈：}Python 3.12+ / FastAPI / LangGraph / Vue 3 / ChromaDB / BM25 / BGE / Git
  \begin{itemize}%
    \item \textbf{Agent 架构：}LangGraph 状态图实现 search--decide--interrupt--execute 工作流，SQLite 持久化 Checkpoint 支持断点续跑
    \item \textbf{混合检索：}BM25 关键词检索（jieba + TF-IDF）+ BGE 向量语义检索，RRF 融合排序兼顾精确匹配与语义理解
    \item \textbf{数据同步：}Markdown + ChromaDB + BM25 + Git 四路实时同步；FastAPI SSE 流式输出
  \end{itemize}}

\cventry{2026.04--至今}{git-reporter：Claude Code 自动化工作报告生成工具}{独立开发者}{}{}{%
  Claude Code Skill 插件，从 Git 提交历史自动生成日报/周报/月报，支持中英双语与参考文档风格学习。\\
  \textbf{技术栈：}PowerShell / Claude Code Skill API / Git
  \begin{itemize}%
    \item 支持多种输出模式，中英文双语自动识别；参考文档学习（RAG）自动模仿已有报告风格
    \item 自动提取代码片段、生成 Mermaid 架构图，一键安装跨平台支持
  \end{itemize}}

\cventry{2024.03--2024.04}{上海因致信息科技有限公司}{前端开发工程师（实习）}{}{}{%
  将 Pisces 项目从 ASP.NET 架构重构至 YTP 框架，负责页面组件化、国际化、权限管理与数据流管理。\\
  \textbf{技术栈：}Vue 2 / Element UI / TypeScript / Axios / i18n
  \begin{itemize}%
    \item 页面重构与组件化开发，实现单页面及主从表 CRUD 功能
    \item 国际化（i18n）支持中英文切换；基于角色的路由与按钮级权限控制
    \item 封装统一 Axios 请求层，统一错误处理与 Token 刷新机制
  \end{itemize}}

\section{技能}
\cvline{Agent 开发}{LangGraph 状态图 Agent、Human-in-the-Loop、持久化 Checkpoint、RAG（BM25 + BGE + RRF 融合）、ChromaDB 向量数据库、SSE 流式交互}
\cvline{LLM 集成}{GLM-4 / ZhipuAI API 调用、Claude Code Skill 开发、结构化 Prompt 设计、多轮对话管理}
\cvline{深度学习}{PyTorch、CNN / ResNet / LSTM / Transformer、MAML 元学习、PCA / t-SNE 降维分析}
\cvline{全栈开发}{后端：Python / FastAPI / SpringBoot；前端：Vue 3 / Vue 2 / Element UI / TypeScript；数据库：MySQL / SQLite}
\cvline{工程工具}{Git 版本控制、LaTeX、MATLAB、AutoDL 云平台、Zotero 文献管理、CET-4}

% \subsection{Vocational}
% \cventry{year--year}{Job title}{Employer}{City}{}{Description}  % arguments 3 to 6 are optional
% \cventry{year--year}{Job title}{Employer}{City}{}{Description}  % arguments 3 to 6 are optional
% \subsection{Miscellaneous}
% \cventry{year--year}{Job title}{Employer}{City}{}{Description line 1\newline{}Description line 2}% arguments 3 to 6 are optional

% \section{Languages}
% \cvlanguage{language 1}{Skill level}{Comment}
% \cvlanguage{language 2}{Skill level}{Comment}
% \cvlanguage{language 3}{Skill level}{Comment}

% \section{Computer skills}
% \cvcomputer{category 1}{XXX, YYY, ZZZ}{category 4}{XXX, YYY, ZZZ}
% \cvcomputer{category 2}{XXX, YYY, ZZZ}{category 5}{XXX, YYY, ZZZ}
% \cvcomputer{category 3}{XXX, YYY, ZZZ}{category 6}{XXX, YYY, ZZZ}

% \section{Interests}
% \cvline{篮球}{\small 体力与技巧}
% \cvline{hobby 2}{\small Description}
% \cvline{hobby 3}{\small Description}

% \renewcommand{\listitemsymbol}{-}  % change the symbol for lists

% \section{Extra 1}
% \cvlistitem{Item 1}
% \cvlistitem{Item 2}
% \cvlistitem[+]{Item 3}  % optional other symbol% XeLaTeX can use any Mac OS X font. See the setromanfont command below.
% Input to XeLaTeX is full Unicode, so Unicode characters can be typed directly into the source.

% The next lines tell TeXShop to typeset with xelatex, and to open and save the source with Unicode encoding.

% !TEX TS-program = xelatex
% !TEX encoding = UTF-8 Unicode

% \section{Extra 2}
% \cvlistdoubleitem[\Neutral]{Item 1}{Item 4}
% \cvlistdoubleitem[\Neutral]{Item 2}{Item 5}
% \cvlistdoubleitem[\Neutral]{Item 3}{}

%% Publications from a BibTeX file
% \nocite{*}
% \bibliographystyle{plain}
% \bibliography{publications}  % 'publications' is the name of a BibTeX file

% \begin{thebibliography}{}
% \bibitem[]{} 移动增强现实可视化综述[C]. ChinaVis 2017.
% \end{thebibliography}


\end{document}


%% end of file `template_en.tex'.

%%% Local Variables:
%%% mode: latex
%%% TeX-command-extra-options: "-shell-escape"
%%% TeX-master: t
%%% TeX-engine: xetex
%%% End:
